\chapter{Independent Increment Processes}

\section{Independent-Increment Processes}

\begin{definition}[Independent Increment Process]\index{independent increment process}
  A stochastic process \(\{X(t)\}_{t \geq 0}\) is said to have \emph{independent increments} if the set of the following random variables,
  \[
    X_{s_1} - X_{0}, X_{t_1} - X_{s_1}, X_{t_2} - X_{s_2}, \dots, X_{t_n} - X_{s_n},
  \]
  are independent for all \(0 \leq s_1 < t_1 \leq s_2 < t_2 \leq \dots \leq s_n < t_n < \infty, n \geq 1\). Furthermore, if the distribution function of \(X_{t + s} - X_{s}\) is independent of \(s\), we say that \(\{X(t)\}_{t \geq 0}\) has \emph{stationary independent increments}. The stochastic process \(\{X(t)\}_{t \geq 0}\) is called stationary independent-increment process or homogeneous independent-increment process if it has stationary independent increments.
\end{definition}

\begin{theorem}
  Suppose that \(X_0 = 0\). For the independent increment process \(\{X(t)\}_{t \geq 0}\), its finite-dimensional distribution can be determined by the distribution function of \(X_t - X_s\) for all \(0 \leq s < t < \infty\).
\end{theorem}

\begin{theorem}
  Suppose that \(\{X(t)\}_{t \geq 0}\) is an independent increment process and \(X_0 = 0\). Then the covariance function \(C_X(s, t) = \Var_X[\min(s, t)]\)
\end{theorem}

\section{Poisson Process}

\begin{definition}[Poisson Process]\index{Poisson process}
  A counting process \(\{N_t\}_{t \geq 0}\) is called a \emph{Poisson process} with rate \(\lambda > 0\) if it satisfies the following conditions:
  \begin{enumerate}
    \item \(N_0 = 0\).
    \item \(\{N_t\}_{t \geq 0}\) has independent increments.
    \item The number of events in any interval of length \(t\) follows a Poisson distribution with mean \(\lambda t\), that is, for all \(s, t \geq 0\) and \(n = 0, 1, \dots\),
    \[
      P\{N_{t + s} - N_s = n\} = \me^{-\lambda t} \frac{(\lambda t)^n}{n!}.
    \]
  \end{enumerate}
\end{definition}

\begin{proposition}
  Let \(\{N_t\}_{t \geq 0}\) be a Poisson process with rate \(\lambda > 0\). Then
  \begin{enumerate}
    \item \(\mu_N(t) = \lambda t\).
    \item \(C_N(s, t) = \Var_N[\min(s, t)] = \lambda \min(s, t)\).
    \item \(R_N(s, t) = \lambda \min(s, t) + \lambda^2 s t\), where
    \[
      \min(s, t) = \frac{s + t - |s - t|}{2}.
    \]
  \end{enumerate}
\end{proposition}

\begin{example}
  Let \(\{N_t\}_{t \geq 0}\) be a Poisson process with rate \(\lambda > 0\). Find the values of
  \begin{enumerate}
    \item \(\E[N_t^2]\),
    \item \(\E[(N_t - N_s)^2]\),
    \item \(P\{N(s) = 1 | N(t) = 4\}\),
    \item \(R_N(s, t)\).
  \end{enumerate}
  for all \(0 \leq s < t < \infty\).
\end{example}
\begin{solution}
  \begin{enumerate}
    \item Since \(N_t\) follows a Poisson distribution with mean \(\lambda t\), we have
    \[
      \E[N_t^2] = \Var_N[N_t] + (\E[N_t])^2 = \lambda t + (\lambda t)^2.
    \]
    \item Since \(N_t - N_s\) follows a Poisson distribution with mean \(\lambda (t - s)\), we have
    \[
      \E[(N_t - N_s)^2] = \Var_N[N_t - N_s] + (\E[N_t - N_s])^2 = \lambda (t - s) + (\lambda (t - s))^2.
    \]
    \item By the independent increment property of the Poisson process, we have
    \begin{align*}
      P\{N(s) &= 1 | N(t) = 4\} = \frac{P\{N(s) = 1\} P\{N(t) - N(s) = 3\}}{P\{N(t) = 4\}} \\
      &= \frac{\me^{-\lambda s} \frac{(\lambda s)^1}{1!} \cdot \me^{-\lambda (t - s)} \frac{(\lambda (t - s))^3}{3!}}{\me^{-\lambda t} \frac{(\lambda t)^4}{4!}} = \frac{4s(t - s)^3}{t^4}.
    \end{align*}
    \item By the definition of the covariance function, we have
    \[
      R_N(s, t) = C_N(s, t) + \mu_N(s) \mu_N(t) = \lambda \min(s, t) + \lambda^2 s t. \qedhere
    \]
  \end{enumerate}
\end{solution}

\section{Gaussian Process}

\begin{definition}[Gaussian Process]\index{Gaussian process}
  A real-valued stochastic process \(\{X_t\}_{t \in T}\) is said to be a \emph{Gaussian process} if all the finite-dimensional distributions have a multivariate normal distribution. That is, for any choice of distinct values \(t_1, t_2, \dots, t_k \in T\), the random vector \(\mathbf{X} = (X_{t_1}, X_{t_2}, \dots, X_{t_k})\) has a multivariate normal distribution with mean vector \(\boldsymbol{\mu} = \E(\mathbf{X})\) and covariance matrix \(\boldsymbol{\Sigma} = \Cov(\mathbf{X}, \mathbf{X})\), which will be denoted as
  \[
    \mathbf{X} \sim \mathrm{N_k}(\boldsymbol{\mu}, \boldsymbol{\Sigma}).
  \]
\end{definition}

\begin{proposition}
  Some properties of Gaussian processes are as follows:
  \begin{enumerate}
    \item Suppose that \(\{X_t\}_{t \in T}\) is a Gaussian process. Then it is strict-sense stationary if and only if it is wide-sense stationary.
    \item \(\{X_t\}_{t \in T}\) is a Gaussian process if and only if any linear combination of its finite-dimensional distributions is normally distributed,i.e., for any \(t_1, t_2, \dots, t_k \in T\) and \(c_1, c_2, \dots, c_k \in \R\), the random variable \(Y = c_1 X_{t_1} + c_2 X_{t_2} + \dots + c_k X_{t_k}\) is normally distributed.
  \end{enumerate}
\end{proposition}

\begin{example}
  Suppose that \(\{X_t\}_{t \in T}\) is a Gaussian process with mean function \(\mu_X(t) = 0\) and autocorrelation function \(R_X(\tau) = \frac{1}{4} \me^{-2|\tau|}\). For all \(t \in T\), find \(P\{0.5 \leq X_t \leq 1\}\).
\end{example}
\begin{solution}
  Since \(\{X_t\}_{t \in T}\) is a Gaussian process with mean function \(\mu_X(t) = 0\) and autocorrelation function \(R_X(\tau) = \frac{1}{4} \me^{-2|\tau|}\), we have
  \[
    \Var_X[X_t] = R_X(0) = \frac{1}{4}.
  \]
  Therefore, \(X_t \sim \mathrm{N}(0, \frac{1}{4})\). Hence,
  \[
    P\{0.5 \leq X_t \leq 1\} = P\left\{\frac{0.5 - 0}{\sqrt{\frac{1}{4}}} \leq Z \leq \frac{1 - 0}{\sqrt{\frac{1}{4}}}\right\} = P\{1 \leq Z \leq 2\} = \Phi(2) - \Phi(1),
  \]
  where \(Z = (X_t - 0)/\sqrt{\frac{1}{4}}\) follows a standard normal distribution.
\end{solution}

\section{Brownian Motion and Wiener Process}

\begin{definition}[Wiener Process]\index{Wiener process}
  A real-valued process \(\{W_t\}_{t \geq 0}\) is called a \emph{Wiener process} or \emph{Brownian motion} with parameter \(\sigma^2\) if it satisfies the following conditions:
  \begin{enumerate}
    \item \(W_0 = 0\).
    \item The increments
    \[
      W_{t_1} - W_{t_0}, W_{t_2} - W_{t_1}, \dots, W_{t_n} - W_{t_{n - 1}}
    \]
    are independent for all \(0 = t_0 < t_1 < t_2 < \dots < t_n\).
    \item \(W_t - W_s\) has a normal distribution with mean \(0\) and variance \(\sigma^2 (t - s)\) for all \(0 \leq s < t\).
  \end{enumerate}
\end{definition}

\begin{proposition}
  Let \(\{W_t\}_{t \geq 0}\) be a Wiener process with parameter \(\sigma^2\). Then
  \begin{enumerate}
    \item \(\mu_W(t) = 0\).
    \item \(C_W(s, t) = R_W(s, t) = \sigma^2 \min(s, t)\).
  \end{enumerate}
\end{proposition}
\begin{proof}
  Since \(W_t\) has a normal distribution with mean \(0\) and variance \(\sigma^2 t\), we have \(\mu_W(t) = 0\). For the covariance function, we have
  \[
    C_W(s, t) = \Cov(W_s, W_t) = \E[W_s W_t] - \E[W_s] \E[W_t] = \E[W_s W_t].
  \]
  \begin{enumerate}
    \item If \(s = t\), then \(C_W(s, t) = \E[W_s^2] = \Var_W[W_s] = \sigma^2 s\).
    \item If \(s < t\), then we can write \(W_t\) as
    \[
      W_t = W_s + (W_t - W_s).
    \]
    Since \(W_s\) and \(W_t - W_s\) are independent and \(W_t - W_s\) has mean \(0\), we have
    \[
      C_W(s, t) = \E[W_s (W_s + (W_t - W_s))] = \E[W_s^2] + \E[W_s (W_t - W_s)] = \sigma^2 s. \qedhere
    \]
    \end{enumerate}
\end{proof}

\begin{proposition}
  The Wiener process \(W_t\) is a Gaussian process.
\end{proposition}